% answers_short.tex —— 2015 年·湖北卷·理科 · 简短答案速查
% 由 tools/make_short_answers.py 自 answers.tex 的 \gksolhead/\gkanswer 提取，
% 并与 2015/2015_数学真题_湖北理_答案要点.md 逐题比对后生成，勿手改（改源 MD 后重跑）。

\begin{gknotice}
\textbf{使用说明}：本节只给最终答案便于快速核对，完整推导见解析版。答案为编写者独立推导，\textbf{非官方标准答案}。
\end{gknotice}

\gkpart{选择题}
\begin{center}\begin{tabular}{ccccccccccc}
\toprule
题号 & 1 & 2 & 3 & 4 & 5 & 6 & 7 & 8 & 9 & 10 \\
\midrule
答案 & A & B & D & C & A & B & B & D & C & B \\
\bottomrule
\end{tabular}\end{center}

\gkpart{填空题}
\begin{itemize}[leftmargin=2.6em]
  \item[\textbf{11.}] \(9\)
  \item[\textbf{12.}] \(2\)
  \item[\textbf{13.}] \(100\sqrt6\)
  \item[\textbf{14.}] (1) \((x-1)^2+(y-\sqrt2)^2=2\)；(2) ①②③
  \item[\textbf{15.}] \(\dfrac12\)
  \item[\textbf{16.}] \(2\sqrt5\)
\end{itemize}

\gkpart{解答题}
\begin{itemize}[leftmargin=2.6em]
  \item[\textbf{17.}] (1) 表格补全为 \(x=\dfrac{\pi}{12},\dfrac{\pi}{3},\dfrac{7\pi}{12},\dfrac{5\pi}{6},\dfrac{13\pi}{12}\)（对应值 \(0,5,0,-5,0\)）；\(f(x)=5\sin\left(2x-\dfrac{\pi}{6}\right)\)；(2) \(\dfrac{\pi}{6}\)
  \item[\textbf{18.}] (1) 当 \(d=2\) 时 \(a_n=2n-1\)，\(b_n=2^{n-1}\)；当 \(d=\dfrac29\) 时 \(a_n=\dfrac{2n+79}{9}\)，\(b_n=9\left(\dfrac29\right)^{n-1}\)；(2) \(T_n=6-\dfrac{2n+3}{2^{n-1}}\)
  \item[\textbf{19.}] (1) 成立（\(PB\perp\) 平面 \(DEF\)）；四面体 \(DBEF\) 是鳖臑，四个面的直角分别为 \(\angle DEB\)、\(\angle DEF\)、\(\angle DFB\)、\(\angle EFB\)；(2) \(\dfrac{\sqrt2}{2}\)
  \item[\textbf{20.}] (1) \(P(Z=8160)=0.3\)，\(P(Z=10200)=0.5\)，\(P(Z=10800)=0.2\)；\(E(Z)=9708\) 元；(2) \(0.973\)
  \item[\textbf{21.}] (1) \(\dfrac{x^2}{16}+\dfrac{y^2}{4}=1\)（即 \(x^2+4y^2=16\)）；(2) 存在最小值，最小值为 \(8\)
  \item[\textbf{22.}] (1) \(f(x)\) 在 \((-\infty,0)\) 上单调递增，在 \((0,+\infty)\) 上单调递减；\(\left(1+\dfrac1n\right)^n<\mathrm{e}\)；(2) \(\dfrac{b_1}{a_1}=2\)，\(\dfrac{b_1b_2}{a_1a_2}=9\)，\(\dfrac{b_1b_2b_3}{a_1a_2a_3}=64\)；推测 \(\dfrac{b_1b_2\cdots b_n}{a_1a_2\cdots a_n}=(n+1)^n\)（证明见解析）；(3) 成立（由均值不等式与 \(\left(1+\dfrac1k\right)^k<\mathrm{e}\) 放缩即得，见解析）
\end{itemize}
